\begin{table}[p]
\centering
\begin{threeparttable}
\caption{Matched comparisons under original and reconstructed evaluation outcomes.}
\label{tab:ftse-outcome-sensitivity}
\scriptsize
\setlength{\tabcolsep}{3pt}
\begin{tabular*}{\textwidth}{@{\extracolsep{\fill}}llcrrr@{}}
\toprule
Comparison & Loss & $\ell$ & $N$ & Original & Reconstructed \\
\midrule
\multicolumn{6}{l}{\textit{Panel A: FTSE 100, full evaluation period}} \\
@@ROW_0@@
@@ROW_1@@
@@ROW_2@@
@@ROW_3@@
@@ROW_4@@
@@ROW_5@@
@@ROW_6@@
@@ROW_7@@
@@ROW_8@@
\addlinespace
\multicolumn{6}{l}{\textit{Panel B: FTSE 100 information comparison in 2020}} \\
@@ROW_9@@
@@ROW_10@@
@@ROW_11@@
@@ROW_12@@
\addlinespace
\multicolumn{6}{l}{\textit{Panel C: DAX information comparison, full evaluation period}} \\
@@ROW_13@@
@@ROW_14@@
@@ROW_15@@
@@ROW_16@@
\bottomrule
\end{tabular*}
\begin{tablenotes}[flushleft]
\footnotesize
\item Notes: Entries are mean loss differences with two-sided descriptive, unadjusted DM $p$-values in brackets. Information is RV-NN-SV minus NN-SV; neural form is RV-NN-SV minus RV-LIN-SV; combination is the criterion-specific combination minus RV-NN-SV. Negative differences favour the first forecast. Both columns use identical forecasts and common eligible origins, retaining the exchange-calendar gap exclusion. Horizons are in trading days; differences are calculated before rounding.
\end{tablenotes}
\end{threeparttable}
\end{table}
